% !PRINTED 333
\weberpage{333}
%<<CONT
that the arc of the curve running between the two curve points $\alpha, f(\alpha)$ and
$\alpha + \delta, f(\alpha + \delta)$ is replaced by the chord. If the second difference is
also taken into account, then the arc of the curve passing through the three successive
points with the abscissae $\alpha, \alpha + \delta,
\alpha + 2\delta$ is replaced by the arc of a parabola
which passes through the same points and whose axis is parallel to the axis of ordinates;
and this parabola will fit the curve still more closely than the chord.
The smaller the interval $\delta$ is, the less influence the higher
differences will have.

How far one has to go in the approximation, then, depends not only on the
accuracy with which one wishes to know $f(x)$, but essentially also
on the density of the values $\alpha, \alpha + \delta, \alpha + 2\delta
\ldots$ for which the function is known. On these theorems rests the
arrangement of our works of tables, especially of the tables of logarithms. There, to be
sure, we are not dealing with integral rational functions; but for
continuous functions in general the same laws hold here. In the tables
one therefore also finds, beside the values $f(\alpha), f(\alpha + \delta),
f(\alpha + 2\delta) \ldots$, the values of the first difference, or of the first two
differences, given. For the customary seven-place tables the
first difference suffices. In the ten-place table ``Thesaurus
logarithmorum'' of \gesperrt{Vega} the second differences are also
given, and must be taken into account in quite precise computations.

Our interpolation formulae can be applied with advantage to compute the roots
of equations, or, more exactly expressed, to improve the approximate values
found by other means.

We can formulate the problem thus: for a given value of $f(x)$ lying between
$f(\alpha)$ and $f(\alpha + \delta)$, the
corresponding value of $x$ is to be found.

We regard $x = \alpha$ as a first approximate value. If we now set
\[
f(x) - f(\alpha) = \varDelta,
\]
then $\varDelta$ is a given value of the same sign as
$\varDelta_\alpha = f(\alpha + \delta) - f(\alpha)$ and absolutely smaller than
$\varDelta_\alpha$. If we further set $x - \alpha = u$, formula (4) gives
\weq{6}{\varDelta = \frac{\varDelta_\alpha u}{\delta} + \frac{\varDelta'_\alpha}{\delta^2}\,\frac{u(u-\delta)}{2} + \cdots}
